\section{Infrastructure}
\label{sec:infra}

\katago{}'s main run used up to 28 V100 GPUs for 19 days~\citep[App.~C]{wu2019accelerating}.
\gofer{} has no compute budget, so v4 is designed to run meaningfully on free resources
first and to use a single rented GPU efficiently later. One codebase and one configuration
format serve four environments: a laptop, a free CI runner, a container on any GPU host,
and a spot instance chosen across clouds. This section describes each, and how the time
and cost of a run are estimated before paying for it.

\subsection{Inference backends}
\label{sec:infra-backends}

The engine evaluates networks in one of two ways. The \emph{in-process} backend links ONNX
Runtime~\citep{onnxruntime2021} into the Go binary through cgo and loads the pinned
shared library (version 1.26.0) at run time. The \emph{sidecar} backend sends batches over
HTTP to a Python inference server. The in-process backend is the default: in the
end-to-end run it finished the same 20-game gate in 1.6\,s instead of 8.9\,s
(\cref{sec:pipeline-eval}). The sidecar needs no C toolchain, which makes it the choice for
quick local tests and for CI runners without cgo. The export step checks PyTorch against
Python ONNX Runtime (\cref{sec:learner-export}). A separate, pre-existing harness checks
that the Go in-process runtime reproduces the Python runtime's outputs. The orchestrator
downloads the pinned ONNX Runtime build itself on Linux and Windows x86-64.

\subsection{Container image}
\label{sec:infra-docker}

A multi-stage Docker~\citep{merkel2014docker} build compiles the engine with the ONNX
Runtime tag in a Go image, then copies the binary into a PyTorch runtime image together
with the pinned ONNX Runtime library, the learner and the orchestrator. The same file
builds a CUDA variant for GPU hosts and a slim CPU variant, selected by a build argument.
A workflow publishes both to the GitHub container registry, which is free for public
repositories, so a rented machine only has to pull an image. The run directory is a
mounted volume. Stopping the container sends \code{SIGTERM}, which the orchestrator handles
cleanly (\cref{sec:orch-resume}). With a restart policy, the container resumes the run
after a crash or reboot.

\subsection{Rented and preemptible GPUs}
\label{sec:infra-cloud}

For cloud GPUs, a SkyPilot~\citep{yang2023skypilot} task file describes the job
independently of the provider. SkyPilot provisions the cheapest acceptable GPU among the
clouds the user has credentials for, preferring spot capacity. Its managed-job mode
relaunches the task after a preemption. The run directory is mounted from an object-storage
bucket, so a replacement machine starts with the state that the preempted one left, and
the orchestrator resumes at the first unfinished stage (\cref{sec:orch-resume}). The
learner's periodic checkpoints (\cref{sec:learner-opt}) bound the training work lost to one
preemption. Two smaller tools cover the remaining cases: a bootstrap script prepares a
plain Linux virtual machine (Go, ONNX Runtime, a CUDA or CPU PyTorch, engine build and
self-test), and an \code{rsync} wrapper copies a run directory to or from any SSH host.
Because shards and generations never change, repeated copies only transfer new files.

\subsection{Free-tier training on GitHub Actions}
\label{sec:infra-actions}

GitHub-hosted runners are free for public repositories~\citep{githubactionsbilling}. A
standard Linux runner has 4 vCPUs and 16\,GB of memory, and a job may run for at most
6 hours~\citep{githubrunners}. v4 uses this as a no-cost training backend:
\begin{enumerate}
  \item A scheduled workflow starts every 6 hours. It is opt-in: it only runs after a
        repository variable is set.
  \item The job restores the run directory from the Actions cache, trains, and saves the
        directory back. The cache evicts entries that have not been accessed for 7
        days~\citep{githubactionscache}; on a 6-hour schedule that does not happen while
        training is enabled.
  \item The orchestrator stops starting new cycles after 4.75 hours. A hard timeout of
        320 minutes stops a long stage early; the per-stage checkpoints of
        \cref{sec:orch-resume} mean the next job continues it.
  \item Each new champion is published as a GitHub Release (\cref{sec:orch-publish}),
        which is permanent. The cache is only working storage.
\end{enumerate}
At most $4 \times 320$ minutes of runner time per day gives about 21 runner-hours, or
about 85 vCPU-hours, per day, with no GPU. We do not claim this is fast. It is a way to
make steady progress on a small network at zero cost, and to keep the whole loop
exercised continuously. The gating configuration for this environment is necessarily weak
(\cref{sec:orch-gate}): with a 192-game cap it can only reliably detect large gains.

\subsection{Estimating time and cost}
\label{sec:infra-estimate}

Every stage writes its duration to an event log. The \code{estimate} command averages
the per-stage times over the most recent trained cycles and projects a run of $C$ cycles:
\begin{equation}
  T(C) = C \sum_{s \in \mathrm{stages}} \bar t_s, \qquad \mathrm{cost}(C) = T(C)\cdot r,
\end{equation}
where $r$ is the hourly price of the instance. With overlap enabled, $\bar t_{\text{self-play}}$
counts only the time the cycle actually waited for self-play, so $T$ is wall-clock time and
not the sum of all work. The intended workflow is to run a few cycles of the target
configuration, read the projection, and then set \code{deadline\_hours} as a hard budget.
An HTML report shows the ladder Elo, each gate's score with its confidence interval, the
time spent per stage in each cycle, and a table of all cycles.

\subsection{Continuous integration}
\label{sec:infra-ci}

A CI workflow runs on every change to the engine, learner, orchestrator, configurations or
infrastructure. It runs the Go shard tests and all Python tests, a real two-cycle pipeline
run on CPU with the smoke configuration, and a build of the CPU container image followed by
a self-test inside it. The workflow uploads the run directory and report as artifacts. At
the time of writing, the CI, free-training and image workflows have been written and
syntax-checked but not yet run on GitHub; the container image has not been built
(\cref{sec:limitations}).
